\scanpage{270}
- \mspage{82} -
\margin{IX}
par un tore. Nous avons posé comme d'habitude \(A_n=A\times_S S_n,\) où \(S_n=\operatorname{Spec}(V/\uline{m}^{n+1}).\) Nous nous intéresserons surtout au cas où \(S\) est un trait, et où la fibre générique \(A_K\) est un schéma abélien à réduction semi-stable sur \(S,\) \(A\) étant le modèle de Néron de \(A_K.\)

Soit \(T_0\) le tore maximal de \(A_0,\) de sorte que \(B_0=A_0^0/T_0\) est un schéma abélien. La construction de 5.1 s'applique alors pour fournir un sous-tore formel \(\hat T\) du complété formel \(\hat A.\) On sait que pour tout \(n,\) le quotient
\[
(7.1.1)\qquad B_n=A_n^0/T_n
\]
est représentable (SGA 3 VI\(_A\) 3.2), plus précisément c'est un schéma plat et de type fini sur \(S_n,\) dont l'image inverse sur \(S_0\) est le schéma abélien \(B_0.\) Par suite \(B_n\) est un \uline{schéma abélien} sur \(S_n,\) et pour \(n\) variable, les \(B_n\) se déduisent les uns des autres par changement de base, donc ils définissent un \uline{schéma abélien formel}
\[
(7.1.2)\qquad\hat B=(B_n)_{n\geq0}.
\]
Nous allons donner des conditions impliquant que ce schéma abélien est \uline{algébrisable} (EGA III 5.4.1), plus précisément que c'est le complété formel d'un schéma abélien projectif sur \(S,\) que nous noterons alors \(B.\) En vertu de loc.cit. \(B\) est déterminé à isomorphisme unique près (comme \(S\)-schéma ; la structure de groupe de \(\hat B\) provient alors de façon unique d'une structure de groupe de \(B,\) en vertu de loc. cit. également). Ce schéma abélien \(B\) sera alors appelé \uline{la partie abélienne} du schéma \(A,\) et parfois noté \(A^{\mathrm{ab}}.\)
\margin{400}
\margin{269}
\margin{400}
